\prefacesection{Abstract}

Market behavior is shaped not only by price movements but also by macroeconomic forces that influence price stability. Current quantitative trading systems built on convolutional neural networks (CNNs) typically represent market dynamics using technical indicator-derived images. However, these models often ignore macroeconomic indicators that significantly influence financial returns. This study introduces TradeLENS (Trade Learning with Economic and Numeric Signals), a CNN-based trading framework that integrates 15 technical indicators and 5 selected macroeconomic indicators into a unified grayscale image representation to determine the appropriate trading action (Buy, Hold, or Sell). For each trading day, a 20$\times$20 pixel image was constructed from a 20-period lookback window of technical and macroeconomic indicators, and labeled as BUY, SELL, or HOLD using a fixed-window turning-point labeling method. Although the proposed model yielded lower classification accuracy (64.76\%) than the technical-only model (69.24\%), it demonstrated superior trading outcomes. Starting from an initial equity of PHP 10,000, it achieved a final equity of PHP 21,715.91 and a mean annual return of 16.55\%, compared to PHP 17,384.65 and 11.32\% for the baseline. SHAP analysis further supported these empirical results by showing that selected macroeconomic indicators received higher feature importance in profitable CNN--MIC predictions. These findings suggest that incorporating macroeconomic indicators, despite not improving classification accuracy, enhances the profitability of the generated trading signals and underscores the practical value of macroeconomic context in CNN-based trading systems.


	%\textit{Keywords}: Deep-learning; image classification; trading indicators; macroeconomic-informed representations; quantitative trading strategies
